\documentclass[aps,prr,twocolumn,groupedaddress]{revtex4-2}

\usepackage{amsmath}
\usepackage{amssymb}
\usepackage{accents}
\usepackage{contour}
\usepackage{relsize}
\usepackage{mathtools}
\usepackage{dsfont}
\usepackage{natbib}
\usepackage[hidelinks]{hyperref}
\usepackage{tikz}
\usepackage[bb=stix,bfbb]{mathalpha}
\usepackage{bm}
\usepackage{graphicx}
\usepackage{placeins}
\usepackage{subcaption}
\usepackage{accents}

\graphicspath{{imgs/}}

\newcommand*{\dt}[1]{%
	\accentset{\mbox{\large\bfseries .}}{#1}}

% Velocity field
\newcommand{\vfield}[1]{
	\boldsymbol{#1}
}
% Pressure field
\newcommand{\pfield}[1]{
	\mathrm{#1}
}
% Inverse
\newcommand{\inv}[1]{
	\mathrm{#1^{-1}}
}

\def\vecU{\boldsymbol{U}}
\def\vecUb{\vfield{U^{0}}}
\def\vecPb{\pfield{P^{0}}}
\def\vecu{\vfield{u}}
\def\vecv{\vfield{v}}
\def\vecw{\vfield{w}}
\def\vecy{\vfield{y}}
\def\vecubar{\bar{\vecu}}
\def\vecvbar{\bar{\vecv}}
\def\vecwbar{\bar{\vecw}}
\def\vecybar{\bar{\vecy}}
\def\bcdot{\boldsymbol{\cdot}}
\def\ie{\textit{i.e.} }
\def\rhs{{r.h.s} }
\def\etc{{etc.} }
\def\Lu{\mathbf{L}}
\def\Nu{\mathbf{N}}
\def\invRec{\inv{\Rey_{c}}}
\def\Rey{\mathrm{Re}}

\def\parv{z}
\def\vecparv{\mathbf{\parv}}


	

\begin{document}

%opening
\title{Asymptotic Center--Manifold for the Navier--Stokes}
\author{Prabal S. Negi}
\email[]{prabal.negi@oist.jp}
%\homepage[]{Your web page}
%\thanks{}
%\altaffiliation{}
\affiliation{Okinawa Institute of Science and Technology Graduate University, Onna, Okinawa 904-0495, Japan}

\date{\today}


\begin{abstract}
	A reduced mathematical model for the flow in an open cavity is presented. The reduction is based on the center manifold theory applied to a perturbation of the original system which allows for a codimension two bifurcation point. The model exhibits many of the key characteristics observed in the flow dynamics including unstable quasi-periodic edge states as well as switching of limit cycles with parameter variations. An explanation for the exchange of stabilities of the limit-cycles is presented based on the cross-coupling terms of the two amplitude equations.
\end{abstract}

\maketitle

\input{intro}
\input{cavity}

\FloatBarrier

\bibliographystyle{apalike}
\bibliography{references}
%\documentclass[aps,prr,twocolumn,groupedaddress]{revtex4-2}

\usepackage{amsmath}
\usepackage{amssymb}
\usepackage{accents}
\usepackage{contour}
\usepackage{relsize}
\usepackage{mathtools}
\usepackage{dsfont}
\usepackage{natbib}
\usepackage[hidelinks]{hyperref}
\usepackage{tikz}
\usepackage[bb=stix,bfbb]{mathalpha}
\usepackage{bm}
\usepackage{graphicx}
\usepackage{placeins}
\usepackage{subcaption}
\usepackage{accents}

\graphicspath{{imgs/}}

\newcommand*{\dt}[1]{%
	\accentset{\mbox{\large\bfseries .}}{#1}}

% Velocity field
\newcommand{\vfield}[1]{
	\boldsymbol{#1}
}
% Pressure field
\newcommand{\pfield}[1]{
	\mathrm{#1}
}
% Inverse
\newcommand{\inv}[1]{
	\mathrm{#1^{-1}}
}

\def\vecU{\boldsymbol{U}}
\def\vecUb{\vfield{U^{0}}}
\def\vecPb{\pfield{P^{0}}}
\def\vecu{\vfield{u}}
\def\vecv{\vfield{v}}
\def\vecw{\vfield{w}}
\def\vecy{\vfield{y}}
\def\vecubar{\bar{\vecu}}
\def\vecvbar{\bar{\vecv}}
\def\vecwbar{\bar{\vecw}}
\def\vecybar{\bar{\vecy}}
\def\bcdot{\boldsymbol{\cdot}}
\def\ie{\textit{i.e.} }
\def\rhs{{r.h.s} }
\def\etc{{etc.} }
\def\Lu{\mathbf{L}}
\def\Nu{\mathbf{N}}
\def\invRec{\inv{\Rey_{c}}}
\def\Rey{\mathrm{Re}}

\def\parv{z}
\def\vecparv{\mathbf{\parv}}


	

\begin{document}

%opening
\title{Asymptotic Center--Manifold for the Navier--Stokes}
\author{Prabal S. Negi}
\email[]{prabal.negi@oist.jp}
%\homepage[]{Your web page}
%\thanks{}
%\altaffiliation{}
\affiliation{Okinawa Institute of Science and Technology Graduate University, Onna, Okinawa 904-0495, Japan}

\date{\today}


\begin{abstract}
	A reduced mathematical model for the flow in an open cavity is presented. The reduction is based on the center manifold theory applied to a perturbation of the original system which allows for a codimension two bifurcation point. The model exhibits many of the key characteristics observed in the flow dynamics including unstable quasi-periodic edge states as well as switching of limit cycles with parameter variations. An explanation for the exchange of stabilities of the limit-cycles is presented based on the cross-coupling terms of the two amplitude equations.
\end{abstract}

\maketitle

\input{intro}
\input{cavity}

\FloatBarrier

\bibliographystyle{apalike}
\bibliography{references}
%\documentclass[aps,prr,twocolumn,groupedaddress]{revtex4-2}

\usepackage{amsmath}
\usepackage{amssymb}
\usepackage{accents}
\usepackage{contour}
\usepackage{relsize}
\usepackage{mathtools}
\usepackage{dsfont}
\usepackage{natbib}
\usepackage[hidelinks]{hyperref}
\usepackage{tikz}
\usepackage[bb=stix,bfbb]{mathalpha}
\usepackage{bm}
\usepackage{graphicx}
\usepackage{placeins}
\usepackage{subcaption}
\usepackage{accents}

\graphicspath{{imgs/}}

\newcommand*{\dt}[1]{%
	\accentset{\mbox{\large\bfseries .}}{#1}}

% Velocity field
\newcommand{\vfield}[1]{
	\boldsymbol{#1}
}
% Pressure field
\newcommand{\pfield}[1]{
	\mathrm{#1}
}
% Inverse
\newcommand{\inv}[1]{
	\mathrm{#1^{-1}}
}

\def\vecU{\boldsymbol{U}}
\def\vecUb{\vfield{U^{0}}}
\def\vecPb{\pfield{P^{0}}}
\def\vecu{\vfield{u}}
\def\vecv{\vfield{v}}
\def\vecw{\vfield{w}}
\def\vecy{\vfield{y}}
\def\vecubar{\bar{\vecu}}
\def\vecvbar{\bar{\vecv}}
\def\vecwbar{\bar{\vecw}}
\def\vecybar{\bar{\vecy}}
\def\bcdot{\boldsymbol{\cdot}}
\def\ie{\textit{i.e.} }
\def\rhs{{r.h.s} }
\def\etc{{etc.} }
\def\Lu{\mathbf{L}}
\def\Nu{\mathbf{N}}
\def\invRec{\inv{\Rey_{c}}}
\def\Rey{\mathrm{Re}}

\def\parv{z}
\def\vecparv{\mathbf{\parv}}


	

\begin{document}

%opening
\title{Asymptotic Center--Manifold for the Navier--Stokes}
\author{Prabal S. Negi}
\email[]{prabal.negi@oist.jp}
%\homepage[]{Your web page}
%\thanks{}
%\altaffiliation{}
\affiliation{Okinawa Institute of Science and Technology Graduate University, Onna, Okinawa 904-0495, Japan}

\date{\today}


\begin{abstract}
	A reduced mathematical model for the flow in an open cavity is presented. The reduction is based on the center manifold theory applied to a perturbation of the original system which allows for a codimension two bifurcation point. The model exhibits many of the key characteristics observed in the flow dynamics including unstable quasi-periodic edge states as well as switching of limit cycles with parameter variations. An explanation for the exchange of stabilities of the limit-cycles is presented based on the cross-coupling terms of the two amplitude equations.
\end{abstract}

\maketitle

\input{intro}
\input{cavity}

\FloatBarrier

\bibliographystyle{apalike}
\bibliography{references}
%\documentclass[aps,prr,twocolumn,groupedaddress]{revtex4-2}

\usepackage{amsmath}
\usepackage{amssymb}
\usepackage{accents}
\usepackage{contour}
\usepackage{relsize}
\usepackage{mathtools}
\usepackage{dsfont}
\usepackage{natbib}
\usepackage[hidelinks]{hyperref}
\usepackage{tikz}
\usepackage[bb=stix,bfbb]{mathalpha}
\usepackage{bm}
\usepackage{graphicx}
\usepackage{placeins}
\usepackage{subcaption}
\usepackage{accents}

\graphicspath{{imgs/}}

\newcommand*{\dt}[1]{%
	\accentset{\mbox{\large\bfseries .}}{#1}}

% Velocity field
\newcommand{\vfield}[1]{
	\boldsymbol{#1}
}
% Pressure field
\newcommand{\pfield}[1]{
	\mathrm{#1}
}
% Inverse
\newcommand{\inv}[1]{
	\mathrm{#1^{-1}}
}

\def\vecU{\boldsymbol{U}}
\def\vecUb{\vfield{U^{0}}}
\def\vecPb{\pfield{P^{0}}}
\def\vecu{\vfield{u}}
\def\vecv{\vfield{v}}
\def\vecw{\vfield{w}}
\def\vecy{\vfield{y}}
\def\vecubar{\bar{\vecu}}
\def\vecvbar{\bar{\vecv}}
\def\vecwbar{\bar{\vecw}}
\def\vecybar{\bar{\vecy}}
\def\bcdot{\boldsymbol{\cdot}}
\def\ie{\textit{i.e.} }
\def\rhs{{r.h.s} }
\def\etc{{etc.} }
\def\Lu{\mathbf{L}}
\def\Nu{\mathbf{N}}
\def\invRec{\inv{\Rey_{c}}}
\def\Rey{\mathrm{Re}}

\def\parv{z}
\def\vecparv{\mathbf{\parv}}


	

\begin{document}

%opening
\title{Asymptotic Center--Manifold for the Navier--Stokes}
\author{Prabal S. Negi}
\email[]{prabal.negi@oist.jp}
%\homepage[]{Your web page}
%\thanks{}
%\altaffiliation{}
\affiliation{Okinawa Institute of Science and Technology Graduate University, Onna, Okinawa 904-0495, Japan}

\date{\today}


\begin{abstract}
	A reduced mathematical model for the flow in an open cavity is presented. The reduction is based on the center manifold theory applied to a perturbation of the original system which allows for a codimension two bifurcation point. The model exhibits many of the key characteristics observed in the flow dynamics including unstable quasi-periodic edge states as well as switching of limit cycles with parameter variations. An explanation for the exchange of stabilities of the limit-cycles is presented based on the cross-coupling terms of the two amplitude equations.
\end{abstract}

\maketitle

\input{intro}
\input{cavity}

\FloatBarrier

\bibliographystyle{apalike}
\bibliography{references}
%\input{main.bbl}

\end{document} 



\end{document} 



\end{document} 



\end{document} 

